% Generated draft: requires mathematical and source review.
\documentclass[11pt]{article}
\usepackage[margin=25mm]{geometry}
\usepackage{fontspec}
\setmainfont{Latin Modern Roman}
\usepackage{amsmath,amssymb,mathtools,unicode-math}
\setmathfont{Latin Modern Math}
\usepackage{enumitem,needspace,xurl,hyperref}
\setlength{\parindent}{0pt}
\setlength{\parskip}{5pt}
\setlength{\emergencystretch}{3em}
\newcommand{\N}{\mathbb{N}}
\newcommand{\Z}{\mathbb{Z}}
\newcommand{\Q}{\mathbb{Q}}
\newcommand{\R}{\mathbb{R}}
\newcommand{\C}{\mathbb{C}}
\newcommand{\F}{\mathbb{F}}
\newcommand{\A}{\mathbb{A}}
\newcommand{\PP}{\mathbb{P}}
\newcommand{\E}{\mathbb{E}}
\newcommand{\eps}{\varepsilon}
\newcommand{\dd}{\,\mathrm{d}}
\begin{document}
\section*{UnsolvedMath ID: 30000295}
\textbf{Source problem code:} catalogue entry
\subsection{Definitions and mathematical statement}
Retain \(\mathcal A^s\) for the solution-approximation class in the source. Let \(u\) solve the intended differential or variational problem with corresponding right-hand side \(f\), and let \(V'\) be the dual space whose norm is intended. The conjectural implication is posed for every exponent \(0<s\leq1/2\) and every such solution \(u\in\mathcal A^s\).

For an admissible partition \(\mathcal T\) of the intended domain \(\Omega\), denote by \(\mathcal P_0(\mathcal T)\) the functions constant on each partition element. Interpreting ``an \(n\)-element partition'' literally, the requested property is
\[
 \exists C>0\ \forall n\in\Z_{>0}\ \exists\mathcal T_n\ \exists f_n\in\mathcal P_0(\mathcal T_n),\qquad
 |\mathcal T_n|=n,\quad\|f-f_n\|_{V'}\leq Cn^{-s}.
\]
The constant may depend on the fixed problem data, \(u\) and \(s\), but must be independent of \(n\). The partition and approximant may depend on \(n\); no nestedness or use of the same partition as the best approximation to \(u\) is asserted. Membership of the piecewise-constant functions in \(V'\) must be part of the missing problem setup. If \(V\) is a normed test space, its dual norm is
\[
 \|g\|_{V'}=\sup_{0\ne v\in V}\frac{|\langle g,v\rangle|}{\|v\|_V}.
\]
The record does not identify \(V\), so a particular Sobolev dual norm or boundary condition is not selected.

A best-approximation interpretation of \(\mathcal A^s\) ordinarily concerns boundedness of a quantity such as \(\sup_n n^s\inf_{v\in\mathcal U_n}\|u-v\|_V\), but neither the discrete classes \(\mathcal U_n\), the solution norm nor the size convention is given. This is explanatory notation, not a definition silently assigned to the source's \(\mathcal A^s\). Similarly, adaptive-mesh conventions may count at most \(n\) elements or \(n\) elements added to an initial mesh, and admissible refinements may not have every exact cardinality. The source needs checking before replacing its literal exact-\(n\) wording by one of these alternatives.

The assertion asks for existence of data approximants with the same decay exponent, not an algorithm producing them or a uniform bound in terms of an explicitly specified class seminorm. The literature triage checked 2026-08-21 reports this implication conjectural in the low-regularity range and says later AFEM work uses data-approximation assumptions without establishing the full implication. Its open label and source-verification report remain dated metadata, not independently verified conclusions.
\subsection{Short English statement}
Does a solution's approximation rate \(n^{-s}\), for \(0<s\leq1/2\), imply the same rate for piecewise-constant approximation of its forcing term in the intended dual norm? The partitions, operator, function spaces and approximation-class definition are missing. The quantifiers require one bound independent of \(n\), with the approximant allowed to vary with \(n\).
\subsection{Sources}
\begin{itemize}
\item Local UnsolvedMath registry, permanent ID 30000295, source code catalogue entry, ``Low-Regularity Right-Hand-Side Approximation in Adaptive Methods,'' including original and cleaned statements and verification metadata.
\item \emph{Mini-Workshop: Convergence of Adaptive Algorithms (2006)}, Oberwolfach Reports, pp.\ 2091--2138, DOI 10.4171/owr/2005/37: \url{https://doi.org/10.4171/owr/2005/37}. The citation uses 2006, while the DOI and proposed-year metadata use 2005.
\item Supplied reference ``Stevenson AFEM paper'': \url{https://doi.org/10.1007/s10208-005-0183-0}.
\item Literature triage checked 2026-08-21; registry record updated 2026-09-16. No external source has been visited or independently checked here.
\end{itemize}
\end{document}
